\documentclass[
	10pt,								% globale Schriftgröße
	parskip=half-,						% setzt Absatzabstand hoch
	paper=a4,							% Format
	english,ngerman,					% lädt Sprachpakete
	]{scrartcl}							% Dokumentenklasse

% //////////////////// Pakete laden ////////////////////
\usepackage{amsmath}			% MUSS vor fontspec geladen werden
\usepackage{mathtools}			% modifiziert amsmath
\usepackage{amssymb}			% mathematische symbole, für \ceckmarks
\usepackage{amsthm}				% für proof
\usepackage{mathrsfs}			% für \mathscr
\usepackage{latexsym}
\usepackage{marvosym}				% für Lightning

\usepackage{fontspec} 			% funktioniert nur mit den neueren Compilern z.B. XeLaTeX
\usepackage{microtype}			% für bessere Worttrennung
\usepackage[ngerman]{babel} 	% Spracheinstellung
\usepackage{lmodern}			% verändert verwendete Schriftart, damit sie weniger pixelig ist

\usepackage{verbatim}
\usepackage{listings}			% Für Quellcode

\usepackage{graphicx}
\usepackage{tabularx}			% für Tabellen mit gleicher Spaltenbreite und automatischen Umbrüchen
\usepackage{fullpage}
\usepackage{multirow}			% für multirow in tabulars
\usepackage{rotate}
\usepackage[cmyk,table]{xcolor} % um Farben zu benutzen, kann mehr als das Paket color
\usepackage[					% Verlinkungen
	colorlinks,					% farbige Schrift, statt farbiger Rahmen
	linktocpage,				% verlinkt im Abb.Verzeichnis Seitenzahl statt Bildunterschrift
	linkcolor=blue				% setzt Farbe der Links auf blau
	]{hyperref}					% nur für digitale Anwendungen, url = "http://www.example.com"
%\usepackage{url}				% für Webadressen wie e-mail usw.: "\url{http://www.example.com}"

\usepackage{enumerate}			% für versch. Aufzählungezeichen wie z.B. a)
\usepackage{xspace}				% folgt ein Leerzeichen nach einem \Befehl, wird es nicht verschluckt.
\usepackage{cancel}				% für das Durchstreichen u.a. in Matheformeln mit \cancel
\usepackage{float}              % zum Forcieren der Position von figure-Umgebungen
\usepackage{ifthen}				% neceassary for \forloop und \whiledo

% zum Zeichnen (u.a. von Graphen)
\usepackage{fp}
\usepackage{tikz}
\usetikzlibrary{tikzmark}			% für \tikzmark{toRemember}
\usetikzlibrary{positioning}	% verbesserte Positionierung der Knoten
\usetikzlibrary{automata}		% für Automaten (GTI)
\usetikzlibrary{arrows}
\usetikzlibrary{shapes}
\usetikzlibrary{decorations.pathmorphing}
\usetikzlibrary{decorations.pathreplacing}
\usetikzlibrary{decorations.shapes}
\usetikzlibrary{decorations.text}

% //////////////////// Syntaxhighlighting ////////////////////
\lstloadlanguages{Python, Haskell, [LaTeX]TeX, Java}
\lstset{
   basicstyle=\footnotesize\ttfamily,	% \scriptsize the size of the fonts that are used for the code
   backgroundcolor = \color{bgcolour},	% legt Farbe der Box fest
   breakatwhitespace=false,	% sets if automatic breaks should only happen at whitespace
   breaklines=true,			% sets automatic line breaking
   captionpos=t,				% sets the caption-position to bottom, t for top
   commentstyle=\color{codeblue}\ttfamily,% comment style
   frame=single,				% adds a frame around the code
   keepspaces=true,			% keeps spaces in text, useful for keeping indentation
							% of code (possibly needs columns=flexible)
   keywordstyle=\bfseries\ttfamily\color{codepurple},% keyword style
   numbers=left,				% where to put the line-numbers;
   							% possible values are (none, left, right)
   numberstyle=\tiny\color{codegreen},	% the style that is used for the line-numbers
   numbersep=5pt,			% how far the line-numbers are from the code
   stepnumber=1,				% nummeriert nur jede i-te Zeile
   showspaces=false,			% show spaces everywhere adding particular underscores;
							% it overrides 'showstringspaces'
   showstringspaces=false,	% underline spaces within strings only
   showtabs=false,			% show tabs within strings adding particular underscores
   flexiblecolumns=false,
   tabsize=1,				% the step between two line-numbers. If 1: each line will be numbered
   stringstyle=\color{orange}\ttfamily,	% string literal style
   numberblanklines=false,				% leere Zeilen werden nicht mitnummeriert
   xleftmargin=1.2em,					% Abstand zum linken Layoutrand
   xrightmargin=0.4em,					% Abstand zum rechten Layoutrand
   aboveskip=2ex, 
}

\lstdefinestyle{py}{
   language=Python,
}
\lstdefinestyle{hs}{
   language=Haskell,
}
\lstdefinestyle{tex}{
	language=[LaTeX]TeX,
	escapeinside={\%*}{*)},     % if you want to add LaTeX within your code
	texcsstyle=*\bfseries\color{blue},% hervorhebung der tex-Schlüsselwörter
	morekeywords={*,$,\{,\},\[,\],lstinputlisting,includegraphics,
	rowcolor,columncolor,listoffigures,lstlistoflistings,
	subsection,subsubsection,textcolor,tableofcontents,colorbox,
	fcolorbox,definecolor,cellcolor,url,linktocpage,subtitle,
	subject,maketitle,usetikzlibrary,node,path,addbibresource,
	printbibliography},% if you want to add more keywords to the set
     numbers=none,
     numbersep=0pt,
     xleftmargin=0.4em,
}

\lstdefinestyle{java}{
	language=Java,
	extendedchars=true,		% lets you use non-ASCII characters;
   						% for 8-bits encodings only, does not work with UTF-8
}

\lstdefinelanguage[x64]{Assembler}     % add a "x64" dialect of Assembler
   [x86masm]{Assembler} % based on the "x86masm" dialect
   % with these extra keywords:
   {morekeywords={CDQE,CQO,CMPSQ,CMPXCHG16B,JRCXZ,LODSQ,MOVSXD, %
                  POPFQ,PUSHFQ,SCASQ,STOSQ,IRETQ,RDTSCP,SWAPGS, %
                  rax,rdx,rcx,rbx,rsi,rdi,rsp,rbp, %
                  r8,r8d,r8w,r8b,r9,r9d,r9w,r9b}
}					% for 8-bits encodings only, does not work with UTF-8

\lstdefinestyle{c}{
	language=c,
	extendedchars=true,		% for 8-bits encodings only, does not work with UTF-8
}

% //////////////////// eigene Kommandos ////////////////////
\newcommand\FU{Freie Universität Berlin\xspace}% benötigt package xspace
\newcommand\gdw{g.\,d.\,w.\xspace}
\newcommand\oBdA{o.\,B.\,d.\,A.\xspace}
\newcommand{\Eu}{\texteuro}
\newcommand\N{\mathbb{N}\xspace}
\newcommand\Q{\mathbb{Q}\xspace}
\newcommand\R{\mathbb{R}\xspace}
\newcommand\Z{\mathbb{Z}\xspace}
\newcommand\ohneNull{\ensuremath{\backslash\lbrace 0\rbrace}}% \{0}
\let\dhALT\dh	% Schreibt Befehl \dh in \dhALT um
\renewcommand\dh{d.\,h.\xspace}	%renew überschreibt command \dh
\newcommand\Bolt{\;\text{\LARGE\raisebox{-0.3em}{\Lightning}\normalsize}\xspace}% Blitz
\newcommand\zz{\ensuremath{\raisebox{+0.25ex}{Z}% zu zeigen
			\kern-0.4em\raisebox{-0.25ex}{Z}%
			\;\xspace}}
\newcommand{\from}{\ensuremath{\colon}}
\newcommand{\floor}[1]{\lfloor{#1}\rfloor}
\newcommand{\ceil}[1]{\lceil{#1}\rceil}
 \renewcommand{\L}{\ensuremath{\mathcal{L}}\xspace}
 \renewcommand{\P}{\ensuremath{\mathcal{P}}\xspace}
 \newcommand{\NL}{\ensuremath{\mathcal{N}\kern-0.2em\mathcal{L}}\xspace}
 \newcommand{\NP}{\ensuremath{\mathcal{NP}}\xspace}

% //////////////////// Mathefunktionen ////////////////////
\DeclareMathOperator{\Landau}{\mathcal{O}}
\DeclareMathOperator{\True}{True}
\DeclareMathOperator{\False}{False}

% //////////////////// eigene Theoreme ////////////////////
\newtheorem{theorem}{Satz}
\newtheorem{corollary}[theorem]{Folgerung}
\newtheorem{lemma}[theorem]{Lemma}
\newtheorem{observation}[theorem]{Beobachtung}
\newtheorem{definition}[theorem]{Definition}
\newtheorem{Literatur}[theorem]{Literatur}
% konfiguriert proof
\makeatletter
\newenvironment{Proof}[1][\proofname]{\par
  \pushQED{\qed}%
  \normalfont \topsep6\p@\@plus6\p@\relax
  \trivlist
  \item[\hskip\labelsep
%         \itshape
        \bfseries
    #1\@addpunct{.}]\ignorespaces
}{%
  \popQED\endtrivlist\@endpefalse
}
\makeatother

% //////////////////// eigene Farben ////////////////////
\let\definecolor=\xdefinecolor
\definecolor{FUgreen}{RGB}{153,204,0}
\definecolor{FUblue}{RGB}{0,51,102}

\definecolor{middlegray}{rgb}{0.5,0.5,0.5}
\definecolor{lightgray}{rgb}{0.8,0.8,0.8}
\definecolor{orange}{rgb}{0.8,0.3,0.3}
\definecolor{azur}{rgb}{0,0.7,1}
\definecolor{yac}{rgb}{0.6,0.6,0.1}
\definecolor{Pink}{rgb}{1,0,0.6}

\definecolor{bgcolour}{rgb}{0.97,0.97,0.97}
\definecolor{codegreen}{rgb}{0,0.6,0}
\definecolor{codegray}{rgb}{0.35,0.35,0.35}
\definecolor{codepurple}{rgb}{0.58,0,0.82}
\definecolor{codeblue}{rgb}{0.4,0.5,1}

% //////////////////// eigene Settings ////////////////////

\textheight = 230mm		% Höhe des Satzspiegels / Layouts
\footskip = 10ex			% Abstand zw. Fußzeile und Grundlinie letzter Textzeile
\parindent 0pt			% verhindert Einrückung der 1. Zeile eines Absatzes
\setkomafont{sectioning}{\rmfamily\bfseries}% setzt Ü-Schriften in Serifen, {disposition}										
\newcommand{\dozent}{Wolfgang Mulzer, Katharina Klost}					% <-- Names des Dozenten eintragen
\newcommand{\tutor}{Tobias Gleißner}						% <-- Name eurer Tutoriun eintragen
\newcommand{\tutoriumNo}{02}				% <-- Nummer im KVV nachschauen
\newcommand{\ubungNo}{01}									% <-- Nummer des Übungszettels
\newcommand{\veranstaltung}{Algorithmen, Datenstrukturen und Datenabstraktion}	% <-- Name der Lehrveranstaltung eintragen
\newcommand{\semester}{Semester}						% <-- z.B. SoSo 17, WiSe 17/18
\newcommand{\studenten}{Nicolas Benjamin, Michael Wernitz}			% <-- Hier eure Namen eintragen
\newcommand{\aufgNo}{3}	

\usepackage{amsmath}
\usepackage{mathtools}
\usepackage{letltxmacro}

% /////////////////////// BEGIN DOKUMENT /////////////////////////
\begin{document}							
% ////////////// Notenspiegel und Logo //////////////
\vspace*{-15ex}							% rückt Logo an den oberen Seitenrand
\makebox[\dimexpr\textwidth+1cm][r]{		%rechtsbündig und geht rechts 1cm über Layout hinaus
	\begin{minipage}{0.5\linewidth}
	\newcounter{AufgNo}
	\setcounter{AufgNo}{\aufgNo}
	\stepcounter{AufgNo}					% AufgNo++
	\newcounter{zahl}
	\def\and{&\xspace}
	\renewcommand{\arraystretch}{1.3}\setlength{\tabcolsep}{1em}
	\begin{tabular}{*{\value{AufgNo}}{|c} |}
		\hline
		\setcounter{zahl}{1}
		\whiledo{\value{zahl} < \value{AufgNo}}{%\AufgNo 
			\thezahl\and\stepcounter{zahl}%
		} $\sum$ \\ \hline
		\setcounter{zahl}{1}
		\whiledo{\value{zahl} < \value{AufgNo}}{%\AufgNo 
			\and\stepcounter{zahl}%
		} \\ \hline
	\end{tabular}

	%\forloop{ct}{1}{\value{ct} < 5}{}
	\end{minipage}
	\hfill
	\begin{minipage}{0.5\linewidth}
	\hfill \ifnum\aufgNo<9%				blendet ab 9 Aufgaben das Logo aus
	\includegraphics[width=0.8\textwidth]{src/fu_logo} % fügt FU-Logo ein
	\fi
	\end{minipage}
}
% ////////////// Daten //////////////
\begin{center}
{\large \dozent}\par
{\huge \veranstaltung, \semester}\par
{\large TutorIn: \tutor, Tutorium \tutoriumNo}\par
{\Large Übung \ubungNo}\par 
{\large \studenten}\par
\today
\end{center}
\vspace{-3ex}							% Abstand
\rule{\linewidth}{0.8pt}				% horizontale Linie

\section{ Rekursionsgleichung \hfill}
{\itshape Lösen Sie die folgenden Rekursionsgleichungen. Beweisen Sie jeweils die Richtigkeit Ihrer Lösung (z.B.: durch vollständige Induktion)}

\begin{flushleft}
	(a)
\end{flushleft}

\begin{center}
$T(1) = 0; T(n) = T(\lfloor{\frac{n}{2}}\rfloor) + T(\lceil{\frac{n}{2}}\rceil) + n$, für n > 1 
\end{center}



\begin{itemize}
\item  Ergebnisse durch Einsetzen von Zweierpotenzen von n\\
	$T(2) = T(1) + T(1) + 2 = 2$ \\
    $T(4) = T(2) + T(2) + 4 = 8$ \\
    $T(8) = T(4) + T(4) + 8 = 24$ \\
    $T(16) = T(8) + T(8) + 16 = 64$ \\
\end{itemize}
 
Daraus erkennt man die geschlossene Formel: $a_{n} = n * log_{2}n$

\begin{proof}
        Es folgt der dazugehörige Beweis durch vollständige Induktion 
        \par
        \begingroup
        \leftskip = 1,35cm
        \noindent
			Dazu betrachten wir n als $n=2^k$\\
			Es wird daher über k induziert
		\par
        \endgroup
        

	\begin{itemize}
	\item Induktionsanfang mit $k=1$, also $2^1$
		$T(2) = T(1) + T(1) + 2$\\
		$T(2) = 2$\\
     	$a_2 = 2 * log(2)$\\
    	$a_2 = 2$\\
     	$=> T(2) = a_2$
        
    \item Induktionsvorraussetzung: $T(n) = a_{n}$\\
    	wobei $T(2^k) = T(\lfloor{\frac{2^{k-1}}{2}}\rfloor) + T(\lceil{\frac{2^{k-1}}{2}}\rceil) + 2^k$\\
        \\
    	und $a_{2^k} = 2^k * log_{2}(2^k)$
    \item Induktionsbehauptung: $a_2^{k+1} = T(2^{k+1})$ 
    	wir gehen dabei davon aus, dass $T(2^{k+1})$ bereits bewiesen wurde, da $T(2^k)$ bzw. $T(n)$ in der Aufgabenstellung gegeben ist.   
    \item Induktionsschritt $k \mapsto k+1$ \\
       
       $T(2^{k+1}) = 2 * T(2^k) + 2^{k+1}$
       \par
       \begingroup
       \leftskip = 1,35cm 
       \noindent
       $ \overset{IV}{=} 2 * 2^k * log_{2}(2^k) + 2^{k+1}$ \\
       \\
       $ = 2^{k+1} * log_{2}(2^k) + 2^{k+1}$ \\
       \\
       $ = 2^{k+1} * (log_{2}(2^k) + 1)$ \\
       \\
       $ = 2^{k+1} * (log_{2}(2^k) + log_{2}2)$\\
       \\
       $ = 2^{k+1} * (log_{2}(2^k * 2^1)$ \\
       \\
       $ = 2^{k+1} * log_{2}(2^{k+1}) \Leftrightarrow a_{k+1}$ 
       \par
       \endgroup  
      
   	 \end{itemize}
	\end{proof}
    
\begin{flushleft}
	(b)
\end{flushleft}
\begin{center}
$S(1) = 1; S(n) = \displaystyle\sum_{i=1}^{n-1}i * S(i)$, für n > 1 
\end{center}

\newpage

\section{Die Ungleichung vom arithmetischen und geometrischen Mittel}


Seien $x_1, x_2, \ldots, x_n \geq 0$. Beweisen Sie
die Ungleichung
\[
P(n): x_1 x_2 \ldots x_n
\leq\left(\frac{x_1+\dots+x_n}{n}\right)^n,
\]
indem Sie folgende Teilschritte bearbeiten:
\begin{enumerate}
  \item $P(2)$ ist eine wahre Aussage.

     \emph{Hinweis}: $(x_1 - x_2)^2 \geq 0$.
  \item
     Wenn $P(n)$ gilt, dann gilt auch $P(n-1)$.

     \emph{Hinweis}: Wie muss man $x_n$ w\"ahlen, damit 
     sich der Wert der rechten Klammer nicht \"andert, wenn
     $x_1, \ldots, x_{n-1}$ schon feststehen?

   \item Aus $P(2)$ und $P(n)$ folgt $P(2n)$.
   \item Folgern Sie nun die Ungleichung. 
\end{enumerate}

\begin{flushleft}
	(1.)
\end{flushleft}

\begin{itemize}
\item
	\begin{proof}
	 P(2) ist wahre Aussage \\
		$$x_1*x_2 \leq (\frac{x_1+x_2}{2})^2$ \\
		$$x_1*x_2 \leq \frac{(x_1+x_2)^2}{4}$  \\
		$$4*x_1*x_2 \leq (x_1+x_2)^2$ \\
		$$4*x_1*x_2 \leq {x_1}^2 + 2*x_1x_2 + {x_2}^2 | -4*x_1x_2$\\
		$$0 \leq {x_1}^2 - 2*x_1x_2 + {x_2}^2 $\\
		$$0 \leq  (x_1-x_2)^2 $ w.A., da $ x_1, x_2 \geq 0 $\\
		
	    
	\end{proof}
\end{itemize}

\newpage



\newpage
\section{Manipulation elementarer Funktionen}

Finden Sie Paare von \"aquivalenten Termen und formen Sie diese
schrittweise ineinander um. Geben Sie die verwendeten Regeln an.
\[
 \log_a\Bigl(n^{\log_b a}\Bigr), \sqrt[b]{\frac{a^n}{a^m}},
 b^{n \log a}, \log_b n, a^{\frac{n-m}{b}}, n(\log a + \log b),
 \log(a^nb^n),a^{(\log b^n)}.
\]
{\itshape Lösungen}
\begin{itemize}
\item $a \Leftrightarrow d$ \\
	  \[
      log_{a}(n^{log_{b}a}) \Leftrightarrow log_{b}a * log_{a}n \Leftrightarrow \frac{log_{a}a}{log_{a}b} * log_{a}n \Leftrightarrow \frac{1}{log_{a}b} * log_{a}n \Leftrightarrow \frac{log_{a}n}{log_{a}b} \Leftrightarrow log_{b}n
      \]
\item $b \Leftrightarrow e$ \\
	  \[
       \sqrt[b]{\frac{a^n}{a^m}} \Leftrightarrow \frac{\sqrt[b]{a^n}}{\sqrt[b]{a^m}} \Leftrightarrow \frac{a^{\frac{n}{b}}}{a^{\frac{m}{b}}} \Leftrightarrow a^{\frac{n-m}{b}}
      \]
\item $c \Leftrightarrow h$ \\
	  \[
        b^{n*log_{2}a} \Leftrightarrow (b^{log_{2}a})^n \Leftrightarrow (2^{log_{2}b * log_{2}a})^n \Leftrightarrow ((2^{log_{2}a})^{log_{2}b})^n \Leftrightarrow (a^{log_{2}b})^n \Leftrightarrow a^{n * log_{2}b}
      \]
\item $ f \Leftrightarrow g$
      \[
       n*(log(a) + log(b)) \Leftrightarrow n * log(a) + n* log(b) \Leftrightarrow log(a^n) + log(b^n) \Leftrightarrow log(a^n*b^n)
      \]
\end{itemize}



% /////////////////////// END DOKUMENT /////////////////////////
\end{document}
